%______________________________________
% WorPT table "[INSERTFILENAME]". It needs \usepackage{worpt} in your preamble (worpt.sty comes with the tables).
% How to put it in your document:
% https://github.com/pmarcum/WorPT/blob/main/latexTableTemplates/[INSERTFILENAME]/readme.md
%______________________________________

%
% Template for WorPT script begins below
% (not for general user unless you are planning to alter the WorPT script!)
%

\newif\ifLandScape                 % establish a boolean variable

% ___________________________________________
%       Default formatting preferences
% ___________________________________________
\LandScapefalse             % landscape mode OFF, unless the table needs it (see MinTitleWidth); \LandScapetrue always uses it

%--- table column widths
% The table is as wide as the page's text. Each team member has a narrow column, just wide enough for its label and the
% numbers at its foot, both printed sideways; the title column gets the width that is left, and a title too long for it
% is cut short with "...". If less than MinTitleWidth is left, the table is set in landscape mode, which is wider.
\def\MinTitleWidth{1.5in}   % narrowest title column before switching to landscape mode
\def\TimelineWidth{0.6ex}   % skinny timeline columns width (one column per month; wider if its labels need it)
\def\TeamCellPadding{1pt}   % room beside the sideways label and numbers in a team member's column

%---- table compactness controls
\def\SpaceBetweenRows{0.8}    % vertical compactness of rows
\def\SpaceBetweenColumns{1pt} % spacing between columns

%---- vertical line colors
\def\TimelineVerticalLineColor{lightgray}       % lines between quarters
\def\TimelineMonthLineColor{lightgray!45}      % lines between months, fainter
\def\TeamVerticalLineColor{lightgray}           % lines between team members

%---- preferences of column headers
\def\TaskTimelineHeaderFontsize{\scriptsize}   % Task Timeline label text size
\def\TeamHeaderFontsize{\scriptsize}           % Team label text size
\def\TotalFteHeaderFontsize{\scriptsize}       % Total FTE label text size (printed sideways)
\def\YearHeaderFontsize{\scriptsize}           % "YEAR1", "YEAR2"... label text size
\def\TaskTitleHeaderFontsize{\normalsize}      % "TASK TITLES" column label text size
\def\YearSliceHeaderFontsize{\scriptsize}      % "1", "2"... quarter label text size (each over its three months)

%--- preferences of task CATEGORY label
\def\TaskCategoryLabel#1#2{{\normalsize\textbf\scshape{#1}}~{\normalsize\textbf{#2}}} % Task Category label format for row item

%--- preferences of TASK label
\def\TaskTitleLabel#1#2{~{\scriptsize{\textbf{\scshape{#1}}}}~{\color{mediumelectricblue}{\footnotesize\textbf{#2}}}} % Task title format for row item

%--- preferences of SUBTASK label: indented under its task
\def\SubtaskIndent{1em}
\def\SubtaskTitleLabel#1#2{\hspace{\SubtaskIndent}\TaskTitleLabel{#1}{#2}} % Subtask title format for row item

%--- preferences of timeline color
\def\TimelineColor{mediumelectricblue}          % Task Timeline cell color (a task, or a subtask)
\def\ParentTimelineColor{mediumelectricblue!35} % Task Timeline cell color of a task that has subtasks (it spans them)

%--- team members' column labels (printed sideways), by institution and funding
\def\FundedUsTeam#1{\small \color{blue}{#1}}    % funded US team member's label
\def\UnfundedUsTeam#1{\small \color{black}{#1}} % unfunded US team member's label
\def\InternationalTeam#1{\small \color{red}{#1}}% international team member's label

%--- the marks in a team member's column: who leads a task, and who works on it (funded, or partly or wholly unfunded)
\def\TeamLeadColor{black}                       % task lead: a solid mark
\def\TeamFundedColor{gray!65}                   % works on the task, funded by this grant: a solid mark
\def\TeamUnfundedColor{gray!80}                 % works on the task, some or all of it not funded by this grant: an open box
\def\TeamMark#1{\TeamSolid{#1}{\dimexpr\linewidth-1pt\relax}}% a solid mark in color #1, as wide as the column allows
\def\TeamOpenMark#1{\TeamOpen{#1}{\dimexpr\linewidth-1pt\relax}}% an open box outlined in color #1, the same size
\def\TeamLeadKey{\TeamSolid{\TeamLeadColor}{1.2em}}      % the marks as the caption shows them
\def\TeamFundedKey{\TeamSolid{\TeamFundedColor}{1.2em}}
\def\TeamUnfundedKey{\TeamOpen{\TeamUnfundedColor}{1.2em}}
\def\TeamOutline{0.6pt}                         % the open box's line
\def\TeamSolid#1#2{\textcolor{#1}{\rule[-0.25ex]{#2}{1.5ex}}}
\def\TeamOpen#1#2{\textcolor{#1}{\fboxsep=0pt\fboxrule=\TeamOutline\fbox{%
  \rule[\dimexpr-0.25ex+\TeamOutline\relax]{0pt}{\dimexpr1.5ex-2\fboxrule\relax}\hspace{\dimexpr#2-2\fboxrule\relax}}}}

%--- preferences for the FTE values: each task's total (right), and each team member's over the whole grant (bottom)
\def\FteTotalFormat#1{\small{#1}}               % Total FTE of a task (or subtask): its own effort
\def\FteRolledUpFormat#1{\small\itshape\color{mediumelectricblue}{#1}} % Total FTE of a task that has subtasks: theirs, added up
\def\FteCategoryFormat#1{\small\bfseries{#1}}  % Total FTE of a task category: its tasks', added up
\def\FteGrandTotalFormat#1{\small\bfseries{#1}}% the whole team's funded and unfunded FTE, at the foot
\def\TeamFundedFteFormat#1{\scriptsize\color{blue}{#1}}% a team member's funded FTE
\def\TeamUnfundedFteFormat#1{\scriptsize{#1}}   % a team member's unfunded FTE
\def\TeamZeroFteFormat#1{\scriptsize\color{lightgray}{#1}}% a team member's FTE when it is zero
\def\TeamFundedFteLabel{\scriptsize\textbf{Funded FTE}}     % label of the bottom row of funded FTE (whole grant)
\def\TeamUnfundedFteLabel{\scriptsize\textbf{Unfunded FTE}} % label of the bottom row of unfunded FTE (whole grant)

%--- the symbols an earlier layout of this table used (kept so captions written for it still work)
\def\TotalFteSumHeaderIcon{\textbf{\large{$\Sigma$}}}
\def\TotalFteUnfundedHeaderIcon{\WorPTnoDollarIcon{-0.4}{0.4mm}{0.2}{0.15}}
\def\TotalFteFundedHeaderIcon{\WorPTdollarIcon{-0.4}{0.2}{0.015}}

%--- Team members' names, for a caption that gives them (the suggested caption does not: the labels are roles, which
% keeps the table anonymous): \RevealIdentities lists them, each as \RevealIdentityFormat{label}{name}
\def\RevealIdentityFormat#1#2{\textbf{#1}: #2}

%--- internal parameters used by script but not for user
\def\LastYearPlusTwo{[INSERTLASTYEARPLUSTWO]}
\def\NumberYears{[INSERTLASTYEAR]}
\def\NumberTeamMinusOne{[INSERTNUMBERTEAMMINUSONE]}
\def\WidestFte{[INSERTWIDESTFTE]}               % the widest number in the Total FTE column
\def\SlicesPerYear{12}       % one timeline column per month (read by WorPT; leave as is)
\def\SlicesPerLabel{3}        % a quarter label spans three months (read by WorPT; leave as is)
\def\QuarterRule{{\color{\TimelineVerticalLineColor}\vrule}}
\def\MonthRule{{\color{\TimelineMonthLineColor}\vrule}}
\def\TeamRule{{\color{\TeamVerticalLineColor}\vrule}}
\def\TimelineSize{\tiny}                        % width control of TASK TIMELINE cells
\def\RevealIdentities{[INSERTIDTONAMELIST]}

%--- the table's columns
\def\[INSERTFILENAME]Columns{
|>{\raggedright\arraybackslash\leavevmode}p{\WorPTtitleWidth} % title column: the width left over
|*{\NumberYears}{@{}%                                % timeline: a column per month; faint lines between months, light ones between quarters
p{\WorPTmonthWidth}@{\MonthRule}p{\WorPTmonthWidth}@{\MonthRule}p{\WorPTmonthWidth}@{\QuarterRule}%
p{\WorPTmonthWidth}@{\MonthRule}p{\WorPTmonthWidth}@{\MonthRule}p{\WorPTmonthWidth}@{\QuarterRule}%
p{\WorPTmonthWidth}@{\MonthRule}p{\WorPTmonthWidth}@{\MonthRule}p{\WorPTmonthWidth}@{\QuarterRule}%
p{\WorPTmonthWidth}@{\MonthRule}p{\WorPTmonthWidth}@{\MonthRule}p{\WorPTmonthWidth}@{}|}%
@{}*{\NumberTeamMinusOne}{>{\centering\arraybackslash\leavevmode}p{\WorPTteamWidth}@{\TeamRule}}% team: a column per member, light lines between
>{\centering\arraybackslash\leavevmode}p{\WorPTteamWidth}@{}%
|>{\centering\arraybackslash\leavevmode}p{\WorPTfteWidth}| % total fte column: as wide as its numbers
% (Written this way for aastex and revtex, which misread a repeat inside a repeat, or a repeat that starts with | followed by @.)
}

%--- the widths: month columns wide enough for the timeline's labels (so the last month of a year is not stretched to
% fit them), team columns wide enough for their numbers and labels, the Total FTE column for its numbers, then the
% rest of the table measured without its titles, and the title column given what is left of the page
\def\[INSERTFILENAME]AtLeast#1#2{\ifdim#1<#2\global#1=#2\fi}
\def\[INSERTFILENAME]Layout{%
  \begingroup
  \csname @floatboxreset\endcsname  % measured as set in a table: some classes (aastex) use a smaller font there
  \renewcommand{\arraystretch}{\SpaceBetweenRows}\setlength{\tabcolsep}{\SpaceBetweenColumns}%
  \global\WorPTmonthWidth=\TimelineWidth\relax
  \setbox0\hbox{\TaskTimelineHeaderFontsize{TIMELINE}}\setbox2\hbox{\TaskTimelineHeaderFontsize{TASK}}%
  \ifdim\wd2>\wd0 \setbox0\box2 \fi
  \[INSERTFILENAME]AtLeast\WorPTmonthWidth{\dimexpr(\wd0+2\tabcolsep-\numexpr12*\NumberYears-1\relax\arrayrulewidth)/\numexpr12*\NumberYears\relax\relax}%
  \setbox0\hbox{\YearHeaderFontsize{YR\NumberYears}}%
  \[INSERTFILENAME]AtLeast\WorPTmonthWidth{\dimexpr(\wd0+2\tabcolsep-11\arrayrulewidth)/12\relax}%
  \setbox0\hbox{\YearSliceHeaderFontsize{4}}%
  \[INSERTFILENAME]AtLeast\WorPTmonthWidth{\dimexpr(\wd0-2\arrayrulewidth)/3\relax}%
  \global\WorPTteamWidth=0pt
  \def\[INSERTFILENAME]Sideways##1{\setbox0\hbox{##1}\[INSERTFILENAME]AtLeast\WorPTteamWidth{\dimexpr\ht0+\dp0+\TeamCellPadding\relax}}%
  \[INSERTFILENAME]Sideways{\FundedUsTeam{Xy}}\[INSERTFILENAME]Sideways{\UnfundedUsTeam{Xy}}\[INSERTFILENAME]Sideways{\InternationalTeam{Xy}}%
  \[INSERTFILENAME]Sideways{\TeamFundedFteFormat{0.}}\[INSERTFILENAME]Sideways{\TeamUnfundedFteFormat{0.}}\[INSERTFILENAME]Sideways{\TeamZeroFteFormat{0.}}%
  \setbox0\hbox{\TeamHeaderFontsize{TEAM}}%
  \[INSERTFILENAME]AtLeast\WorPTteamWidth{\dimexpr(\wd0+2\tabcolsep-\NumberTeamMinusOne\arrayrulewidth)/\numexpr\NumberTeamMinusOne+1\relax\relax}%
  \global\WorPTfteWidth=0pt
  \setbox0\hbox{\FteTotalFormat{\WidestFte}}\[INSERTFILENAME]AtLeast\WorPTfteWidth{\wd0}%
  \setbox0\hbox{\FteRolledUpFormat{\WidestFte}}\[INSERTFILENAME]AtLeast\WorPTfteWidth{\wd0}%
  \setbox0\hbox{\FteCategoryFormat{\WidestFte}}\[INSERTFILENAME]AtLeast\WorPTfteWidth{\wd0}%
  \setbox0\hbox{\FteGrandTotalFormat{\WidestFte}}\[INSERTFILENAME]AtLeast\WorPTfteWidth{\wd0}%
  \setbox0\hbox{\TotalFteHeaderFontsize{TOTAL FTE}}\[INSERTFILENAME]AtLeast\WorPTfteWidth{\dimexpr\ht0+\dp0\relax}% (lies sideways)
  \global\WorPTtitleWidth=0pt \WorPTmeasuringtrue
  \sbox\WorPTmeasureBox{\expandafter\tabular\expandafter{\[INSERTFILENAME]Columns}\[INSERTFILENAME]Body\endtabular}%
  \ifLandScape\global\WorPTlandscapetrue\else\global\WorPTlandscapefalse\fi
  \WorPTstraySpace{\[INSERTFILENAME]Columns}\setbox\WorPTmeasureBox\hbox to\dimexpr\wd\WorPTmeasureBox-\WorPTstray\relax{}%
  \global\WorPTtitleWidth=\dimexpr\textwidth-\wd\WorPTmeasureBox\relax
  \ifdim\WorPTtitleWidth<\MinTitleWidth\relax\global\WorPTlandscapetrue\fi
  \ifWorPTlandscape\global\WorPTtitleWidth=\dimexpr\textheight-\wd\WorPTmeasureBox\relax\fi
  \ifdim\WorPTtitleWidth<\MinTitleWidth\relax
    \global\WorPTtitleWidth=\MinTitleWidth\relax
    \PackageWarningNoLine{worpt}{The schedule table is wider than the page even in landscape mode}%
  \fi
  \endgroup}

\newenvironment{[INSERTFILENAME]}{%
\[INSERTFILENAME]Layout
\ifWorPTlandscape
  \begin{sidewaystable*}
\else
  \begin{table*}[h!]     % \begin{table} needed to make bottom table caption have same width as table
\fi
\renewcommand{\arraystretch}{\SpaceBetweenRows} % controls vertical row separation/spacing
\setlength{\tabcolsep}{\SpaceBetweenColumns}    % controls horizontal column separation/spacing
\begin{WorPTtabular}{\[INSERTFILENAME]Columns}\[INSERTFILENAME]Body\end{WorPTtabular}}
{%
\ifWorPTlandscape
  \end{sidewaystable*}
\else
  \end{table*}
\fi
}

%--- the table itself: measured once without its titles, then set
\def\[INSERTFILENAME]Body{%

% ------------------------- First line of headers
\cline{2-[INSERTNUMBERYEARSLICESPLUSONE]}
\multicolumn{1}{c|}{} &                                                  % title column
\multicolumn{[INSERTNUMBERYEARSLICES]}{c|}{\TaskTimelineHeaderFontsize{TASK}} & % schedule columns
\multicolumn{[INSERTNUMBERTEAM]}{c}{} &                                  % team columns
\multicolumn{1}{c}{}\\                                                   % total fte column

% ------------------------- Second line of headers
\multicolumn{1}{c|}{} &                                                    % title column
\multicolumn{[INSERTNUMBERYEARSLICES]}{c|}{\TaskTimelineHeaderFontsize{TIMELINE}} & % schedule columns
\multicolumn{[INSERTNUMBERTEAM]}{c}{} &                                    % team columns
\multicolumn{1}{c}{}\\                                                     % total fte column
\cline{2-[INSERTLASTCOLUMN]}

% -------------------------  Third line of headers
\multicolumn{1}{c|}{}                                                 % title column
[INSERTYEARLABELLIST] &                                               % year label list (YR1, YR2, ...)
\multicolumn{[INSERTNUMBERTEAM]}{c|}{\TeamHeaderFontsize{TEAM}} &    % team columns
\multicolumn{1}{c|}{}\\                                              % total fte column
\cline{1-1}

% COMMENT: below template generates a single year label.  The script glues them together in a single string
% COMMENT: and replaces the INSERTYEARLABELLIST above.
[STARTYEARLABELTEMPLATE]
 & \multicolumn{[INSERTNUMBERSLICESPERYEAR]}{c|}{\YearHeaderFontsize{YR[INSERTYEARNUMBER]}}
[ENDYEARLABELTEMPLATE]

% -------------------------  Fourth line of headers
\multicolumn{1}{|>{\leavevmode}p{\WorPTtitleWidth}|}{\WorPTunlessMeasuring{\TaskTitleHeaderFontsize{TASK TITLES}}} % title column
[INSERTYEARSLICELABELLIST]   % quarter labels (1,2,3,4 1,2,3,4...), each over its months
[INSERTTEAMLABELLIST] &      % team members' labels, sideways
\rotatebox{90}{\TotalFteHeaderFontsize{TOTAL FTE}}\\ % total fte column
\hline

% COMMENT: below template generates a single quarter label (1 through 4), spanning its three month columns, with the line
% COMMENT: that follows it (INSERTLABELRULE: a quarter line, or the year line at the end of a year).
% COMMENT: The script glues them together into single string and replaces the above INSERTYEARSLICELABELLIST
[STARTYEARSLICELABELTEMPLATE]
 & \multicolumn{[INSERTSLICESPERLABEL]}{@{}c[INSERTLABELRULE]}{\YearSliceHeaderFontsize{[INSERTSLICENUMBER]}}
[ENDYEARSLICELABELTEMPLATE]

% COMMENT: below templates make a team member's column label, sideways, by the member's institution and funding.
% COMMENT: The script glues them together in team order and replaces INSERTTEAMLABELLIST above.
[STARTFUNDEDUSTEAMLABELTEMPLATE]
 & \rotatebox{90}{\FundedUsTeam{[TEAMID]}}
[ENDFUNDEDUSTEAMLABELTEMPLATE]
[STARTUNFUNDEDUSTEAMLABELTEMPLATE]
 & \rotatebox{90}{\UnfundedUsTeam{[TEAMID]}}
[ENDUNFUNDEDUSTEAMLABELTEMPLATE]
[STARTINTERNATIONALTEAMLABELTEMPLATE]
 & \rotatebox{90}{\InternationalTeam{[TEAMID]}}
[ENDINTERNATIONALTEAMLABELTEMPLATE]

% begin main contents of table
[INSERTLISTOFCATEGORYSECTIONS]

% ------------------------- the team members' FTE over the whole grant, funded and unfunded
\hline
\multicolumn{[INSERTNUMBERYEARSLICESPLUSONE]}{|r|}{\WorPTunlessMeasuring{\TeamFundedFteLabel}}
[INSERTTEAMFUNDEDFTELIST] &      % each member's funded FTE
\FteGrandTotalFormat{[INSERTFTEFUNDEDTOTAL]}\\
\multicolumn{[INSERTNUMBERYEARSLICESPLUSONE]}{|r|}{\WorPTunlessMeasuring{\TeamUnfundedFteLabel}}
[INSERTTEAMUNFUNDEDFTELIST] &    % each member's unfunded FTE
\FteGrandTotalFormat{[INSERTFTEUNFUNDEDTOTAL]}\\
\hline

% COMMENT: below templates make one team member's FTE at the foot of the table, sideways: funded, unfunded, or zero.
% COMMENT: The script glues them together in team order and replaces INSERTTEAMFUNDEDFTELIST and INSERTTEAMUNFUNDEDFTELIST.
[STARTTEAMFUNDEDFTETEMPLATE]
 & \rotatebox{90}{\TeamFundedFteFormat{[INSERTFTE]}}
[ENDTEAMFUNDEDFTETEMPLATE]
[STARTTEAMUNFUNDEDFTETEMPLATE]
 & \rotatebox{90}{\TeamUnfundedFteFormat{[INSERTFTE]}}
[ENDTEAMUNFUNDEDFTETEMPLATE]
[STARTTEAMZEROFTETEMPLATE]
 & \rotatebox{90}{\TeamZeroFteFormat{[INSERTFTE]}}
[ENDTEAMZEROFTETEMPLATE]

% COMMENT ****************************** CATEGORY TEMPLATE ***************************************
% COMMENT: the below produces the row for a task category label
% COMMENT: The category sections are put together as a single string delimited by end of line character and then
% COMMENT: replace the INSERTLISTOFCATEGORYSECTIONS above, constructing the full table
[STARTCATEGORYSECTIONTEMPLATE]
\WorPTfitTitle\TaskCategoryLabel{[INSERTCATEGORYLABEL]}{[INSERTCATEGORYTITLE]} % category identification
[INSERTEMPTYYEARSLICESLIST] &  % empty year slices
[INSERTEMPTYTEAMCELLSLIST]     % empty team cells
\FteCategoryFormat{[INSERTCATEGORYFTE]}\\ % the category's total fte
[INSERTTASKSUNDERTHISCATEGORY] % tasks under this category
\hline
[ENDCATEGORYSECTIONTEMPLATE]

% COMMENT: the below template makes a single empty year slice. The script figures out how many to make and
% COMMENT: then glues them together into a single string and replaces the INSERTEMPTYYEARSLICESLIST above
[STARTUNCOLOREDTIMELINETEMPLATE]
 & \TimelineSize{}
[ENDUNCOLOREDTIMELINETEMPLATE]

% COMMENT ****************************** ROW TEMPLATES ***************************************
% COMMENT: the below produce a single row for a task under a category: a task, a task that has subtasks (its
% COMMENT: timeline spans theirs and its FTE adds theirs up), or a subtask. The script then makes several such rows
% COMMENT: and glues them toegether, delimited by a end of line character, and this single string replaces the
% COMMENT: INSERTTASKSUNDERTHISCATEGORY above in the CATEGORYSECTIONTEMPLATE above.
[STARTTASKROWTEMPLATE]
% begin task
\WorPTfitTitle\TaskTitleLabel{[INSERTTASKLABEL]}{[INSERTTASKTITLE]} % title of this task
[INSERTYEARSLICESLIST] &  % timeline for this task
[INSERTTEAMCELLSFORTHISTASK]  % who leads it and who works on it
\FteTotalFormat{[INSERTFTETOTAL]}\\
[ENDTASKROWTEMPLATE]
[STARTPARENTTASKROWTEMPLATE]
% begin task
\WorPTfitTitle\TaskTitleLabel{[INSERTTASKLABEL]}{[INSERTTASKTITLE]} % title of this task
[INSERTYEARSLICESLIST] &  % timeline for this task, spanning its subtasks'
[INSERTTEAMCELLSFORTHISTASK]  % (its subtasks have the people)
\FteRolledUpFormat{[INSERTFTETOTAL]}\\
[ENDPARENTTASKROWTEMPLATE]
[STARTSUBTASKROWTEMPLATE]
% begin task
\WorPTfitTitle\SubtaskTitleLabel{[INSERTTASKLABEL]}{[INSERTTASKTITLE]} % title of this subtask
[INSERTYEARSLICESLIST] &  % timeline for this subtask
[INSERTTEAMCELLSFORTHISTASK]  % who leads it and who works on it
\FteTotalFormat{[INSERTFTETOTAL]}\\
[ENDSUBTASKROWTEMPLATE]

% COMMENT: the below template makes a single year slice that falls within the task timeline, helping to construct
% COMMENT: a colored task timeline. The script picks this or the UNCOLORED template above the ROW TEMPLATE written above,
% COMMENT: and puts them together into a single string that replaces INSERTYEARSLICESLIST above
[STARTCOLOREDTIMELINETEMPLATE]
 & \TimelineSize\WorPTcellcolor{\TimelineColor}
[ENDCOLOREDTIMELINETEMPLATE]
[STARTPARENTCOLOREDTIMELINETEMPLATE]
 & \TimelineSize\WorPTcellcolor{\ParentTimelineColor}
[ENDPARENTCOLOREDTIMELINETEMPLATE]

% COMMENT: the below templates make one team member's cell in a task's row: the task lead, a member who works on the
% COMMENT: task funded by this grant, one whose work on it is partly or wholly unfunded, or one not on the task.
% COMMENT: The script glues them together in team order and replaces INSERTTEAMCELLSFORTHISTASK (and, empty, the
% COMMENT: INSERTEMPTYTEAMCELLSLIST of a category's row). Each ends with the & that starts the next cell.
[STARTTEAMLEADCELLTEMPLATE]
\TeamMark{\TeamLeadColor} &
[ENDTEAMLEADCELLTEMPLATE]
[STARTTEAMFUNDEDCELLTEMPLATE]
\TeamMark{\TeamFundedColor} &
[ENDTEAMFUNDEDCELLTEMPLATE]
[STARTTEAMUNFUNDEDCELLTEMPLATE]
\TeamOpenMark{\TeamUnfundedColor} &
[ENDTEAMUNFUNDEDCELLTEMPLATE]
[STARTTEAMEMPTYCELLTEMPLATE]
{} &
[ENDTEAMEMPTYCELLTEMPLATE]

}

% COMMENT: The caption for this table requires variables (just as the isANONtravel did). The below template defines
% COMMENT: the variables used in the caption formulation that is given in the github help page. Specifically,
% COMMENT: the below defines the names that go with the abbreviations used in the table.  IF YOU LEAVE OUT
% COMMENT: THIS INFORMATION FROM THE CAPTIOn, this table could be used as an anonymous one.
% COMMENT: the below is the template used for each item in the list. The items are delimited by a comma
[STARTIDTONAMETEMPLATE]
\RevealIdentityFormat{[INSERTTEAMID]}{[INSERTTEAMNAME]}
[ENDIDTONAMETEMPLATE]
